% Options for packages loaded elsewhere
% Options for packages loaded elsewhere
\PassOptionsToPackage{unicode}{hyperref}
\PassOptionsToPackage{hyphens}{url}
\PassOptionsToPackage{dvipsnames,svgnames,x11names}{xcolor}
%
\documentclass[
  letterpaper,
  DIV=11,
  numbers=noendperiod]{scrartcl}
\usepackage{xcolor}
\usepackage{amsmath,amssymb}
\setcounter{secnumdepth}{5}
\usepackage{iftex}
\ifPDFTeX
  \usepackage[T1]{fontenc}
  \usepackage[utf8]{inputenc}
  \usepackage{textcomp} % provide euro and other symbols
\else % if luatex or xetex
  \usepackage{unicode-math} % this also loads fontspec
  \defaultfontfeatures{Scale=MatchLowercase}
  \defaultfontfeatures[\rmfamily]{Ligatures=TeX,Scale=1}
\fi
\usepackage{lmodern}
\ifPDFTeX\else
  % xetex/luatex font selection
\fi
% Use upquote if available, for straight quotes in verbatim environments
\IfFileExists{upquote.sty}{\usepackage{upquote}}{}
\IfFileExists{microtype.sty}{% use microtype if available
  \usepackage[]{microtype}
  \UseMicrotypeSet[protrusion]{basicmath} % disable protrusion for tt fonts
}{}
\makeatletter
\@ifundefined{KOMAClassName}{% if non-KOMA class
  \IfFileExists{parskip.sty}{%
    \usepackage{parskip}
  }{% else
    \setlength{\parindent}{0pt}
    \setlength{\parskip}{6pt plus 2pt minus 1pt}}
}{% if KOMA class
  \KOMAoptions{parskip=half}}
\makeatother
% Make \paragraph and \subparagraph free-standing
\makeatletter
\ifx\paragraph\undefined\else
  \let\oldparagraph\paragraph
  \renewcommand{\paragraph}{
    \@ifstar
      \xxxParagraphStar
      \xxxParagraphNoStar
  }
  \newcommand{\xxxParagraphStar}[1]{\oldparagraph*{#1}\mbox{}}
  \newcommand{\xxxParagraphNoStar}[1]{\oldparagraph{#1}\mbox{}}
\fi
\ifx\subparagraph\undefined\else
  \let\oldsubparagraph\subparagraph
  \renewcommand{\subparagraph}{
    \@ifstar
      \xxxSubParagraphStar
      \xxxSubParagraphNoStar
  }
  \newcommand{\xxxSubParagraphStar}[1]{\oldsubparagraph*{#1}\mbox{}}
  \newcommand{\xxxSubParagraphNoStar}[1]{\oldsubparagraph{#1}\mbox{}}
\fi
\makeatother


\usepackage{longtable,booktabs,array}
\usepackage{calc} % for calculating minipage widths
% Correct order of tables after \paragraph or \subparagraph
\usepackage{etoolbox}
\makeatletter
\patchcmd\longtable{\par}{\if@noskipsec\mbox{}\fi\par}{}{}
\makeatother
% Allow footnotes in longtable head/foot
\IfFileExists{footnotehyper.sty}{\usepackage{footnotehyper}}{\usepackage{footnote}}
\makesavenoteenv{longtable}
\usepackage{graphicx}
\makeatletter
\newsavebox\pandoc@box
\newcommand*\pandocbounded[1]{% scales image to fit in text height/width
  \sbox\pandoc@box{#1}%
  \Gscale@div\@tempa{\textheight}{\dimexpr\ht\pandoc@box+\dp\pandoc@box\relax}%
  \Gscale@div\@tempb{\linewidth}{\wd\pandoc@box}%
  \ifdim\@tempb\p@<\@tempa\p@\let\@tempa\@tempb\fi% select the smaller of both
  \ifdim\@tempa\p@<\p@\scalebox{\@tempa}{\usebox\pandoc@box}%
  \else\usebox{\pandoc@box}%
  \fi%
}
% Set default figure placement to htbp
\def\fps@figure{htbp}
\makeatother





\setlength{\emergencystretch}{3em} % prevent overfull lines

\providecommand{\tightlist}{%
  \setlength{\itemsep}{0pt}\setlength{\parskip}{0pt}}



 


\KOMAoption{captions}{tableheading}
\makeatletter
\@ifpackageloaded{caption}{}{\usepackage{caption}}
\AtBeginDocument{%
\ifdefined\contentsname
  \renewcommand*\contentsname{Table of contents}
\else
  \newcommand\contentsname{Table of contents}
\fi
\ifdefined\listfigurename
  \renewcommand*\listfigurename{List of Figures}
\else
  \newcommand\listfigurename{List of Figures}
\fi
\ifdefined\listtablename
  \renewcommand*\listtablename{List of Tables}
\else
  \newcommand\listtablename{List of Tables}
\fi
\ifdefined\figurename
  \renewcommand*\figurename{Figure}
\else
  \newcommand\figurename{Figure}
\fi
\ifdefined\tablename
  \renewcommand*\tablename{Table}
\else
  \newcommand\tablename{Table}
\fi
}
\@ifpackageloaded{float}{}{\usepackage{float}}
\floatstyle{ruled}
\@ifundefined{c@chapter}{\newfloat{codelisting}{h}{lop}}{\newfloat{codelisting}{h}{lop}[chapter]}
\floatname{codelisting}{Listing}
\newcommand*\listoflistings{\listof{codelisting}{List of Listings}}
\makeatother
\makeatletter
\makeatother
\makeatletter
\@ifpackageloaded{caption}{}{\usepackage{caption}}
\@ifpackageloaded{subcaption}{}{\usepackage{subcaption}}
\makeatother
\usepackage{bookmark}
\IfFileExists{xurl.sty}{\usepackage{xurl}}{} % add URL line breaks if available
\urlstyle{same}
\hypersetup{
  pdftitle={Assignment 1: Overview of Visualizations},
  pdfauthor={Mykola Utkin},
  colorlinks=true,
  linkcolor={blue},
  filecolor={Maroon},
  citecolor={Blue},
  urlcolor={Blue},
  pdfcreator={LaTeX via pandoc}}


\title{Assignment 1: Overview of Visualizations}
\usepackage{etoolbox}
\makeatletter
\providecommand{\subtitle}[1]{% add subtitle to \maketitle
  \apptocmd{\@title}{\par {\large #1 \par}}{}{}
}
\makeatother
\subtitle{Principles of Data Visualization}
\author{Mykola Utkin}
\date{}
\begin{document}
\maketitle

\renewcommand*\contentsname{Table of contents}
{
\hypersetup{linkcolor=}
\setcounter{tocdepth}{3}
\tableofcontents
}

\section{Visualization 1 - Grouped Bar
Chart}\label{visualization-1---grouped-bar-chart}

\textbf{Image:}

\begin{figure}[H]

{\centering \pandocbounded{\includegraphics[keepaspectratio]{Assignment_1_2_rand_viz_files/mediabag/Screenshot-2024-12-1.png}}

}

\caption{Data Science Salaries by Role}

\end{figure}%

\textbf{Source:}
\url{https://www.kdnuggets.com/data-science-salaries-job-market-analysis-2024-2025}

\textbf{Chart type:} Horizontal Grouped Bar Chart (salary ranges by job
role and experience level)

\subsection{Mapping Analysis}\label{mapping-analysis}

\begin{longtable}[]{@{}
  >{\raggedright\arraybackslash}p{(\linewidth - 4\tabcolsep) * \real{0.4318}}
  >{\raggedright\arraybackslash}p{(\linewidth - 4\tabcolsep) * \real{0.2500}}
  >{\raggedright\arraybackslash}p{(\linewidth - 4\tabcolsep) * \real{0.3182}}@{}}
\toprule\noalign{}
\begin{minipage}[b]{\linewidth}\raggedright
What's being shown
\end{minipage} & \begin{minipage}[b]{\linewidth}\raggedright
Data type
\end{minipage} & \begin{minipage}[b]{\linewidth}\raggedright
Channel used
\end{minipage} \\
\midrule\noalign{}
\endhead
\bottomrule\noalign{}
\endlastfoot
Job roles (Data Scientist, ML Engineer, Data Analyst, etc.) & Nominal &
Position (Y-axis) \\
Experience levels (Entry, Mid, Senior) & Ordinal & Color hue, grouped
bars \\
Salary amount (USD) & Quantitative & Bar length (position on X-axis) \\
Salary labels & Quantitative & Text annotations on bars \\
\end{longtable}

\subsection{Effectiveness Note}\label{effectiveness-note}

The grouped horizontal bars allow for clear comparison between
experience levels within each role. The color differentiation for levels
is effective, but the overall readability could improve with consistent
ordering of levels and spacing between role groups to avoid visual
crowding.

\begin{center}\rule{0.5\linewidth}{0.5pt}\end{center}

\section{Visualization 2 - Multi-line
Chart}\label{visualization-2---multi-line-chart}

\textbf{Image:}

\begin{figure}[H]

{\centering \pandocbounded{\includegraphics[keepaspectratio]{Assignment_1_2_rand_viz_files/mediabag/Screenshot-2024-12-11.png}}

}

\caption{Salary Trends by Experience Level}

\end{figure}%

\textbf{Source:}
\url{https://www.kdnuggets.com/data-science-salaries-job-market-analysis-2024-2025}

\textbf{Chart type:} Multi-line Chart (Average salary over time)

\subsection{Mapping Analysis}\label{mapping-analysis-1}

\begin{longtable}[]{@{}lll@{}}
\toprule\noalign{}
What's being shown & Data type & Channel used \\
\midrule\noalign{}
\endhead
\bottomrule\noalign{}
\endlastfoot
Year (2020--2024) & Temporal & Position (X-axis) \\
Average salary (USD) & Quantitative & Position (Y-axis), line height \\
Experience level & Ordinal & Color hue, line style \\
Salary trend & Quantitative & Line slope/angle \\
\end{longtable}

\subsection{Effectiveness Note}\label{effectiveness-note-1}

The separate color-coded lines clearly display progression across
experience levels. The viewer easily sees that senior and executive
salaries increase faster than entry-level ones. Adding data labels or
smoothing could improve the precision of interpretation without clutter.

\begin{center}\rule{0.5\linewidth}{0.5pt}\end{center}

\section{Visualization 3 - Heatmap}\label{visualization-3---heatmap}

\textbf{Image:}

\begin{figure}[H]

{\centering \pandocbounded{\includegraphics[keepaspectratio]{"D:/RProjects/Data_vizualisation/data-science-salarie.png"}}

}

\caption{Salary Heatmap by Role and Experience}

\end{figure}%

\textbf{Source:}
\url{https://www.refontelearning.com/salary-guide/data-science-salaries-unveiled}

\textbf{Chart type:} Heatmap (Matrix showing salaries by role and
experience)

\subsection{Mapping Analysis}\label{mapping-analysis-2}

\begin{longtable}[]{@{}lll@{}}
\toprule\noalign{}
What's being shown & Data type & Channel used \\
\midrule\noalign{}
\endhead
\bottomrule\noalign{}
\endlastfoot
Job roles & Nominal & Position (Y-axis) \\
Experience levels & Ordinal & Position (X-axis) \\
Salary (USD) & Quantitative & Color intensity (dark = higher salary) \\
Numeric value & Quantitative & Text label (optional) \\
\end{longtable}

\subsection{Effectiveness Note}\label{effectiveness-note-2}

The sequential color scale effectively highlights the salary gradient
across roles and experience levels. Darker tones draw attention to
top-paying roles, such as Data Architect and Chief Data Officer. A color
legend or numeric scale would further improve interpretability.

\begin{center}\rule{0.5\linewidth}{0.5pt}\end{center}

\section{Visualization 4 - Stacked Bar
Chart}\label{visualization-4---stacked-bar-chart}

\textbf{Image:}

\begin{figure}[H]

{\centering \pandocbounded{\includegraphics[keepaspectratio]{Assignment_1_2_rand_viz_files/mediabag/__results___38_0.png}}

}

\caption{Distribution of Data Science Jobs}

\end{figure}%

\textbf{Source:}
\url{https://www.kaggle.com/code/mohamedmagdy191/data-science-salaries-eda-visualization}

\textbf{Chart type:} Horizontal Stacked Bar Chart (100\% normalized)

\subsection{Mapping Analysis}\label{mapping-analysis-3}

\begin{longtable}[]{@{}
  >{\raggedright\arraybackslash}p{(\linewidth - 4\tabcolsep) * \real{0.4318}}
  >{\raggedright\arraybackslash}p{(\linewidth - 4\tabcolsep) * \real{0.2500}}
  >{\raggedright\arraybackslash}p{(\linewidth - 4\tabcolsep) * \real{0.3182}}@{}}
\toprule\noalign{}
\begin{minipage}[b]{\linewidth}\raggedright
What's being shown
\end{minipage} & \begin{minipage}[b]{\linewidth}\raggedright
Data type
\end{minipage} & \begin{minipage}[b]{\linewidth}\raggedright
Channel used
\end{minipage} \\
\midrule\noalign{}
\endhead
\bottomrule\noalign{}
\endlastfoot
Job titles (8--10 roles) & Nominal & Position (Y-axis) \\
Experience level (EN, MI, SE, EX) & Ordinal & Bar segments, color hue \\
Share of experience levels & Quantitative & Segment length (width) \\
Total composition & Quantitative & Full bar = 100\% \\
Years 2020-2023 & Temporal & Position (x-axis) \\
\end{longtable}

\subsection{Effectiveness Note}\label{effectiveness-note-3}

The 100\% stacked format shows workforce composition across roles
effectively. Roles such as Data Scientist have balanced experience
distribution, while ML Research Scientist is dominated by
senior/executive positions. Color-coding is intuitive, though adding
labels would make small segments easier to interpret.

\begin{center}\rule{0.5\linewidth}{0.5pt}\end{center}

\section{Visualization 5 - Pie Chart}\label{visualization-5---pie-chart}

\textbf{Image:}

\begin{figure}[H]

{\centering \pandocbounded{\includegraphics[keepaspectratio]{Assignment_1_2_rand_viz_files/mediabag/__results___24_0.png}}

}

\caption{Distribution of Job Titles in Data Science}

\end{figure}%

\textbf{Source:}
\url{https://www.kaggle.com/code/mohamedmagdy191/data-science-salaries-eda-visualization}

\textbf{Chart type:} Pie Chart with percentage labels

\subsection{Mapping Analysis}\label{mapping-analysis-4}

\begin{longtable}[]{@{}
  >{\raggedright\arraybackslash}p{(\linewidth - 4\tabcolsep) * \real{0.4318}}
  >{\raggedright\arraybackslash}p{(\linewidth - 4\tabcolsep) * \real{0.2500}}
  >{\raggedright\arraybackslash}p{(\linewidth - 4\tabcolsep) * \real{0.3182}}@{}}
\toprule\noalign{}
\begin{minipage}[b]{\linewidth}\raggedright
What's being shown
\end{minipage} & \begin{minipage}[b]{\linewidth}\raggedright
Data type
\end{minipage} & \begin{minipage}[b]{\linewidth}\raggedright
Channel used
\end{minipage} \\
\midrule\noalign{}
\endhead
\bottomrule\noalign{}
\endlastfoot
Job titles (Data Engineer, Data Scientist, etc.) & Nominal &
Slice/sector position \\
Percentage of total jobs & Quantitative & Slice angle and area total of
100\% \\
Percentage labels & Quantitative & Text values at the middle of the
sector \\
Categories & Nominal & Color hue \\
Categories & Categorical & From largest percentage to least in non
clockwise order \\
\end{longtable}

\subsection{Effectiveness Note}\label{effectiveness-note-4}

The pie chart immediately communicates that Data Engineer and Data
Scientist dominate the market, comprising roughly two-thirds of all
roles. The descending slice sizes are easy to read. Still, small slices
are difficult to compare --- a sorted horizontal bar chart could provide
a more accurate visual ranking.

\begin{center}\rule{0.5\linewidth}{0.5pt}\end{center}




\end{document}
